\chapter{子模块性能仿真}

特别说明：在本章节中，为确保仿真结果输出的可靠性与图表展示的清晰度，除了典型标称状态（Nominal/TT, $3.3\,\mathrm{V}$, $25^\circ\mathrm{C}$）外，我们特意选取了以下四个最具代表性的极限工艺角与 PVT 条件下的仿真波形进行重叠绘制与分析。选择这四个工艺角作为典型边界条件的理由如下：

\begin{enumerate}
    \item \textbf{Asym\_FS (Fast NMOS, Slow PMOS, $3.3\,\mathrm{V}$, $25^\circ\mathrm{C}$):} 
    该不对称工艺角会导致 NMOS 和 PMOS 的载流子迁移率与跨导（$g_m$）发生严重错配。对于全差分架构（FDA/PGA）和电流镜偏置电路，这种错配会从物理底层破坏电路的对称性。选取该状态是评估共模反馈（CMFB）稳定性、共模抑制比（CMRR）恶化程度以及系统直流失调（Offset）漂移的最佳“不对称极限”场景。

    \item \textbf{Asym\_SF (Slow NMOS, Fast PMOS, $3.3\,\mathrm{V}$, $25^\circ\mathrm{C}$):} 
    与 Asym\_FS 呈相反的错配极端。在具有级联结构的运算放大器中，该工艺角极易导致不同支路的静态工作点发生偏移，进而挤压晶体管的输出电压裕度（$V_{ov}$）。将其纳入典型绘制，旨在双向验证电路在极端工艺偏斜下，是否仍能维持合理的静态偏置与大信号动态范围。

    \item \textbf{BC (Best Case - FF, $3.6\,\mathrm{V}$, $-40^\circ\mathrm{C}$):} 
    在该状态下，器件速度最快、供电电压最高且温度最低，晶体管的跨导和运放的增益带宽积（GBW）将飙升至系统最大值。选取该工艺角主要用于严苛考核闭环系统的\textbf{高频稳定性}，验证在 GBW 暴增的情况下，非主极点是否会侵蚀相位裕度（Phase Margin），以及瞬态阶跃响应中是否会暴露出严重的过冲（Overshoot）或高频振荡（Ringing）风险。

    \item \textbf{WCLowtemp (Worst Case Low Temp - SS, $3.0\,\mathrm{V}$, $-40^\circ\mathrm{C}$):} 
    该状态叠加了最慢的工艺、极低的供电电压以及低温下的阈值电压（$V_{th}$）上升，导致晶体管驱动能力降至冰点。这是评估系统\textbf{“速度与建立瓶颈”}的绝对恶劣条件。选取该角用于验证在带宽与压摆率（Slew Rate）被严重压缩的极限下，系统是否依然能够满足 $700\,\mathrm{kHz}$ 的目标带宽，以及能否在规定的时钟周期内完成 14-bit 精度的建立（Settling）。
\end{enumerate}
完整的PVT与工艺角仿真可见仿真文件中。

% -------------------------------
% 跨阻放大器（TIA）单模块仿真
% @ 模块基本测试条件：
%   1. 电源3.3V
%   2. 输入电流范围0.1n~1uA
% @ 模块技术细节：
%   1. 采用两级ota架构运放作为主运放
%   2. 外围采用Rf=100k，Cf=1p（可以更改Cf以找到最佳的值）
% @ 模块设计图见图二
% -------------------------------
\section{跨阻放大器（TIA）单模块仿真}

\subsection{输入共模范围与直流响应}

\textbf{仿真设置：}TIA 核心采用两级 OTA 架构，闭环反馈网络配置为 $R_f=100\,\mathrm{k}\Omega$、$C_f=1\,\mathrm{pF}$。在 $3.3\,\mathrm{V}$ 典型供电电压下，对输入直流光电流进行参数扫描（DC Sweep），扫描范围设定为 $0.1\,\mathrm{nA}$ 至 $1\,\mathrm{\mu A}$，重点评估 $50\,\mathrm{nA}$ 典型工作点附近的直流特性。仿真所用 testbench 如图\ref{fig:tia_icmr_tb}所示。
\begin{figure}[H]
    \centering
    \includegraphics[width=0.3\textwidth]{tia_icmr_tb.png}
    \caption{TIA输入共模范围与直流扫描仿真 testbench}
    \label{fig:tia_icmr_tb}
\end{figure}

\textbf{仿真输入：}直流光信号等效电流源 $I_{photo}$。

\textbf{仿真输出：}TIA 核心运放反相输入端（虚拟地）电平与输出节点直流电压轨迹。输出节点电压随输入直流电流变化的特性曲线，在四个极端工艺角（Asym\_FS, Asym\_SF, BC, WCLowtemp）下的重叠对比图如图\ref{fig:tia_icmr_curve}所示。
\begin{figure}[H]
    \centering
    \includegraphics[width=0.8\textwidth]{tia_icmr_curve.png}
    \caption{TIA输出直流电平随输入光电流变化的扫描曲线}
    \label{fig:tia_icmr_curve}
\end{figure}

\textbf{结果分析：} 仿真曲线表明，在伪差分/单端跨阻拓扑下，随着输入光电流的增加，TIA 输出节点的直流电平（即共模电压基线）必然发生线性偏移。在稳定输出共模目标为 $V_{cm} = 1.65\,\mathrm{V}$ 的前提下，当输入电流达到 $330\,\mathrm{nA}$ 时，输出电平已产生约 $0.03\,\mathrm{V}$ 的直流偏差。对于 14-bit 高精度系统而言，前端如此量级的共模漂移将严重挤压后级 ADC 的动态量程，属不可容忍的系统级误差。各工艺角下偏差趋势保持一致，说明此系拓扑固有特性而非工艺失配引起。

\textbf{解决方案：} 针对 TIA 输出端固有的共模漂移问题，本设计在系统链路层进行校正：信号自 TIA 输出后，将级联至全差分单端转差分模块（Single-to-Differential Amplifier），并在该级及后续的 PGA 模块中引入强健的共模反馈（CMFB）环路，以强制重置并稳定全链路的中心共模电平至绝对的 $1.65\,\mathrm{V}$。

\subsection{跨阻增益，带宽与频率响应}

\textbf{仿真设置：}基于 $50\,\mathrm{nA}$ 的直流偏置工作点，注入幅度为 $1\,\mathrm{nA}$ 的小信号交流电流（AC current）。交流扫频（AC Sweep）范围设定为 $1\,\mathrm{Hz}$ 至 $100\,\mathrm{MHz}$。为探究反馈电容对系统高频特性的影响并寻找最佳相位补偿点，对 $C_f$ 进行了参数化扫描（$0.5\,\mathrm{pF}, 1\,\mathrm{pF}, 2\,\mathrm{pF}$）。仿真所用 testbench 同上，如图\ref{fig:tia_icmr_tb}所示。

\textbf{仿真输出：}TIA 输出节点的闭环跨阻增益（以 $\mathrm{dB\Omega}$ 计）与相位响应曲线。在四个极端工艺角下的幅相频特性如图\ref{fig:tia_ac_curve}所示。

\begin{figure}[H]
    \centering
    \includegraphics[width=0.8\linewidth]{tia_ac_curve.png}
    \caption{TIA 闭环跨阻增益与相位频率响应仿真结果}
    \label{fig:tia_ac_curve}
\end{figure}

提取幅频特性的关键指标，不同补偿电容下的结果如表\ref{tab:tia_ac}所示。
\begin{table}[H]
\centering
\caption{不同反馈电容 $C_f$ 下的 TIA 频响特性汇总}
\label{tab:tia_ac}
\begin{tabular}{cccc}
\toprule
$C_f$ 容值 & 低频跨阻增益 & $-3\,\mathrm{dB}$ 闭环带宽 & 频响平坦度/稳定性评估 \\
\midrule
$0.5\,\mathrm{pF}$ & $100\,\mathrm{dB\Omega}$ & $5.5\,\mathrm{MHz}$ & 带宽过大，伴随轻微 Peaking，裕度偏低 \\
$1\,\mathrm{pF}$   & $100\,\mathrm{dB\Omega}$ & $1.8\,\mathrm{MHz}$ & 带宽适中，最平坦巴特沃斯响应，裕度优良 \\
$2\,\mathrm{pF}$   & $100\,\mathrm{dB\Omega}$ & $870\,\mathrm{kHz}$ & 带宽受限（接近系统要求），过度阻尼 \\
\bottomrule
\end{tabular}
\end{table}

\textbf{结果分析：} TIA 在闭环状态下，严格意义的环路相位裕度（Phase Margin）需通过 STB 分析测量，但通过观测闭环幅频响应的峰化（Peaking）现象可直观判断系统的阻尼状态。仿真显示：较小的 $C_f$（如 $0.5\,\mathrm{pF}$）虽能拓展带宽，但会导致高频段增益凸起，时域表现为瞬态振铃；较大的 $C_f$（如 $2\,\mathrm{pF}$）能提供绝对的稳定性，但牺牲了宝贵的带宽裕量。综合权衡噪声滤波与后级 $700\,\mathrm{kHz}$ 级联收缩需求，选定 $1\,\mathrm{pF}$ 作为最佳相位补偿点，此时无明显峰化且具有充足的安全裕量。

\subsection{噪声分析}
\textbf{仿真设置：}运行 Noise 仿真，输出探针置于 TIA 输出节点，输入参考源设定为光电流源接口。积分频率范围覆盖系统完整带宽（$1\,\mathrm{Hz} \sim 1\,\mathrm{MHz}$）。仿真所用 testbench 同上，如图\ref{fig:tia_icmr_tb}所示。

\textbf{仿真输出：}折算至 TIA 差分输出节点的电压噪声谱密度（Noise Spectral Density, NSD）与频段内累积的积分噪声（Integrated Noise）曲线。在四个极端工艺角下的统计学波形如图\ref{fig:tia_noise_curve}所示。
\begin{figure}[H]
    \centering
    \includegraphics[width=0.5\linewidth]{tia_noise_curve}
    \caption{TIA 等效输出噪声谱密度与累积积分噪声曲线}
    \label{fig:tia_noise_curve}
\end{figure}

\textbf{仿真结果：} 结合曲线与 Calculator 函数提取，噪声性能指标汇总如表\ref{tab:tia_noise}所示。

\begin{table}[H]
\centering
\caption{TIA 核心噪声性能仿真结果}
\label{tab:tia_noise}
\begin{tabular}{ccc}
\toprule
评估指标 & 关键评估带宽/频点 & 仿真量测值 \\
\midrule
等效输入电压噪声密度 & $100\,\mathrm{kHz}$ & $49\,\mathrm{nV/\sqrt{Hz}}$ \\
输出端总积分噪声 & 积分上限 $700\,\mathrm{kHz}$ & $\approx 50\,\mathrm{\mu V_{rms}}$ \\
\bottomrule
\end{tabular}
\end{table}

\textbf{结果分析：} TIA 作为模拟前端的第一级“关口”，其噪声底色将直接被后续的 PGA 级联放大，是系统 SNR 最敏感的环节。仿真证实，通过增大 OTA 输入对管尺寸及提升尾电流偏置，有效压低了闪烁噪声拐角及热噪声平台。即使在 $700\,\mathrm{kHz}$ 宽频积分下，总积分噪声仍控制在 $\approx 50\,\mathrm{\mu V_{rms}}$ 量级，接近 14-bit ADC ($3\,\mathrm{V_{pp}}$) 约 $53\,\mathrm{\mu V_{rms}}$ 的理想量化噪声下限，满足系统的苛刻底噪预算。

\subsection{瞬态大信号阶跃响应}
\textbf{仿真设置：}进行长时 Transient 仿真，在系统趋于稳态后，输入光电流施加从 $50\,\mathrm{nA}$ 突跳至 $60\,\mathrm{nA}$ 的陡峭阶跃激励，重点观察 TIA 输出电压的大信号压摆与小信号建立行为。

\textbf{仿真输出：}TIA 输出节点电压随时间的瞬态响应波形。四个不同工艺角下的阶跃恢复与建立曲线如图\ref{fig:tia_step_curve}所示。
\begin{figure}[H]
    \centering
    \includegraphics[width=0.8\linewidth]{tia_step_curve.png}
    \caption{TIA 大信号阶跃瞬态响应及建立时间细节}
    \label{fig:tia_step_curve}
\end{figure}

\textbf{仿真结果：} 从阶跃沿触发至输出电压稳定在最终目标电平 $\pm 0.1\%$ 误差带内的响应参数如表\ref{tab:tia_step}所示：

\begin{table}[H]
\centering
\caption{不同工艺角下 TIA 瞬态阶跃建立时间}
\label{tab:tia_step}
\begin{tabular}{ccc}
\toprule
仿真工艺角 (Corners) & $0.1\%$ 建立时间 ($t_s$) & 最终稳态建立电平 \\
\midrule
Asym\_FS & $750\,\mathrm{ns}$ & $1.65538\,\mathrm{V}$ \\
Asym\_SF & $750\,\mathrm{ns}$ & $1.65538\,\mathrm{V}$ \\
BC (Best Case) & $750\,\mathrm{ns}$ & $1.65538\,\mathrm{V}$ \\
WCLowtemp (Worst Case) & $750\,\mathrm{ns}$ & $1.65539\,\mathrm{V}$ \\
\bottomrule
\end{tabular}
\end{table}

\textbf{结果分析：} 瞬态阶跃响应直观地反映了频率域的主极点带宽与相位裕度。在所有四个严苛的工艺角测试中，输出信号均表现出极佳的时域干净度，无明显的过冲（Overshoot）或振铃（Ringing）现象，证实了 $1\,\mathrm{pF}$ 密勒补偿电容的绝对稳定性。即使在最慢的 WCLowtemp 条件下，建立时间亦保持在 $750\,\mathrm{ns}$ 水平，末端稳态电平的微小直流偏移处于 $\mu\mathrm{V}$ 级别，完全在系统容错规范内，保证了高频大信号冲击下系统的动态保真度。

\subsection{实际指标总结}
基于前文全面的交直流及瞬态仿真，TIA 单模块的核心物理级性能参数最终收敛并汇总如下：

\begin{table}[H]
\centering
\caption{TIA 模块最终仿真核心指标汇总清单}
\label{tab:tia_final_specs}
\begin{tabular}{lp{6cm}p{6cm}}
\toprule
\textbf{设计参数指标} & \textbf{仿真验证结果} & \textbf{设计目标与备注} \\
\midrule
有效输入电流范围 & $0.1\,\mathrm{nA} \sim 330\,\mathrm{nA}$ (未饱和前) & 目标 $0.1\,\mathrm{nA} \sim 1\,\mathrm{\mu A}$（后半段由于共模偏移需后级校正） \\
低频跨阻增益 ($Z_T$) & $100\,\mathrm{dB\Omega}$ ($100\,\mathrm{k}\Omega$) & 完美契合 $R_f = 100\,\mathrm{k}\Omega$ 设计值 \\
闭环 $-3\,\mathrm{dB}$ 带宽 & $1.8\,\mathrm{MHz}$ (@ $C_f = 1\,\mathrm{pF}$) & 远超 $700\,\mathrm{kHz}$，为级联带宽收缩留足裕量 \\
输出等效积分噪声 & $\approx 50\,\mathrm{\mu V_{rms}}$ ($0 \sim 700\,\mathrm{kHz}$) & 满足 14-bit 高精度 ADC 极低底噪需求 \\
大信号阶跃建立时间 & $750\,\mathrm{ns}$ ($0.1\%$ 精度误差带) & 时域平滑无振铃，压摆率充足 \\
\bottomrule
\end{tabular}
\end{table}

% -------------------------------
% 全差分放大器（FDA）仿真
% @ 模块基本测试条件：
%   1. 电源3.3V
%   2. 标准共模输出1.65V
% @ 模块技术细节：
%   1. 采用两级架构，第一级为Cascode，第二级为普通的共源极，提高输出摆幅
% @ 模块设计图见图三
% -------------------------------

\section{全差分放大器（FDA）单模块仿真}

\subsection{输入共模范围与直流稳态特性}

\textbf{仿真设置：}在典型供电电压 $3.3\,\mathrm{V}$ 条件下，设置输出共模反馈（CMFB）的目标基准电压为 $1.65\,\mathrm{V}$。对输入端共模电平进行直流扫描（DC Sweep），扫描区间设定为 $0.5\,\mathrm{V}$ 至 $2.5\,\mathrm{V}$。
仿真所用 testbench 如图\ref{fig:fda_icmr_tb}所示。

\begin{figure}[H]
    \centering
    \includegraphics[width=0.4\linewidth]{fda_icmr_tb.png}
    \caption{FDA输入共模范围仿真testbench}
    \label{fig:fda_icmr_tb}
\end{figure}

\textbf{仿真输出：} 提取 FDA 核心晶体管的工作区状态（Region），并绘制输出共模电压随输入共模电平变化的传输特性曲线。
在四个极限工艺角及标称状态下，输出共模电压的扫描曲线对比如图\ref{fig:fda_icmr_curve}所示。

\begin{figure}[H]
    \centering
    \includegraphics[width=0.7\linewidth]{fda_icmr_curve.png}
    \caption{FDA输入共模范围仿真结果}
    \label{fig:fda_icmr_curve}
\end{figure}

\textbf{仿真结果：} 从曲线拐点及管子工作状态提取出的有效输入共模范围如下表\ref{tab:fda_icmr}所示。

\begin{table}[H]
\centering
\caption{FDA 有效输入共模范围 (ICMR) 仿真结果}
\label{tab:fda_icmr}
\begin{tabular}{ccc}
\toprule
\textbf{工艺角 (Corner)} & \textbf{共模下限 ($V_{\mathrm{icmr,min}}$)} & \textbf{共模上限 ($V_{\mathrm{icmr,max}}$)} \\
\midrule
Nominal ($25^\circ\mathrm{C}$, $3.3\mathrm{V}$) & $0.40\,\mathrm{V}$ & $2.20\,\mathrm{V}$ \\
BC (Best Case) & $0.35\,\mathrm{V}$ & $2.30\,\mathrm{V}$ \\
WCLowtemp (Worst Case) & $0.45\,\mathrm{V}$ & $2.10\,\mathrm{V}$ \\
\bottomrule
\end{tabular}
\end{table}

\textbf{结果分析：} 其判定基准为：在该共模区间内，FDA 内部的输入差分对管、尾电流源以及折叠/共源共栅（Cascode）负载管均深处饱和区（Saturation Region），且 CMFB 环路具备足够的增益将输出共模精准锁定在 $1.65\,\mathrm{V}$ 附近。若前级（如 TIA）输出的共模电平漂移超出此 $0.4\,\mathrm{V} \sim 2.2\,\mathrm{V}$ 的安全裕度，将导致 FDA 增益骤降并引发严重的非线性失真。因此，在系统级联时需严格做好级间的共模电平匹配。

\subsection{开环增益，带宽与相位裕度}

\textbf{仿真设置：}切断外部反馈网络，将 FDA 置于开环测试模式（或使用大电感大电容隔离交直流环路）。输入共模偏置在 $1.65\,\mathrm{V}$，在输入端注入单位交流小信号（AC=1）。
仿真所用 testbench 如图\ref{fig:fda_openloop_tb}所示。

\begin{figure}[H]
    \centering
    \includegraphics[width=0.4\linewidth]{fda_openloop_tb.png}
    \caption{FDA开环增益，带宽与相位裕度仿真testbench}
    \label{fig:fda_openloop_tb}
\end{figure}

\textbf{仿真输出：} FDA 差分输出节点的 Bode 图（幅频与相频响应）。
在四个不同工艺角下的开环频率响应曲线如图\ref{fig:fda_openloop_curve}所示。

\begin{figure}[H]
    \centering
  \includegraphics[width=0.8\linewidth]{fda_openloop_curve.png}
  \caption{FDA开环增益，带宽与相位裕度仿真结果}
  \label{fig:fda_openloop_curve}
\end{figure}

\textbf{仿真结果：} 提取典型状态下的开环频响特征参数，如表\ref{tab:fda_openloop}所示。

\begin{table}[H]
\centering
\caption{FDA 开环频率响应核心指标}
\label{tab:fda_openloop}
\begin{tabular}{cccc}
\toprule
\textbf{测试状态} & \textbf{低频直流增益 ($A_{v0}$)} & \textbf{增益带宽积 (GBW)} & \textbf{相位裕度 (PM)} \\
\midrule
典型工艺角 & $138\,\mathrm{dB}$ & $130\,\mathrm{MHz}$ & $\approx 63^\circ$ \\
\bottomrule
\end{tabular}
\end{table}

\textbf{结果分析：} $138\,\mathrm{dB}$ 的超高开环增益是 14-bit 系统的坚实保障，它能将闭环工作时的静态稳态误差（$1/A\beta$）压缩至极其微小的水平，确保高精度信号的无损传递。约 $63^\circ$ 的相位裕度是理想的巴特沃斯（Butterworth）响应点，预示着闭环阶跃响应将具有极小的过冲。然而，$130\,\mathrm{MHz}$ 的 GBW 在面对较大闭环增益配置时，可用带宽将被显著压缩，要求在系统设计时必须严格进行增益-带宽的折中分配。

\subsection{闭环稳定性（STB）仿真}

\textbf{仿真设置：}将 FDA 置于闭环工作模式。利用 \texttt{iprobe} 探针打断反馈主环路进行 STB（稳定性）分析。为全面评估系统鲁棒性，通过改变反馈电阻网络，扫描闭环增益从 $0\,\mathrm{dB}$ 到 $60\,\mathrm{dB}$，步长设定为 $20\,\mathrm{dB}$。
仿真所用 testbench 如图\ref{fig:fda_stb_tb}所示。

\begin{figure}[H]
    \centering
    \includegraphics[width=0.4\linewidth]{fda_stb_tb.png}
    \caption{FDA闭环稳定性（stb）仿真testbench}
    \label{fig:fda_stb_tb}
\end{figure}

\textbf{仿真输出：} 不同闭环增益配置下的环路增益（Loop Gain）与相位响应曲线。
在四个不同工艺角下的响应曲线重叠图如图\ref{fig:fda_stb_curve}所示。

\begin{figure}[H]
    \centering
  \includegraphics[width=0.8\linewidth]{fda_stb_curve.png}
  \caption{FDA闭环稳定性（stb）仿真结果}
  \label{fig:fda_stb_curve}
\end{figure}

\textbf{仿真结果：} 随闭环增益增加，环路带宽发生收缩，具体数值分布如表\ref{tab:fda_stb}所示：

\begin{table}[H]
\centering
\caption{不同闭环增益配置下的环路稳定性仿真结果}
\label{tab:fda_stb}
\begin{tabular}{cccc}
\toprule
\textbf{闭环增益设定} & \textbf{闭环 $-3\,\mathrm{dB}$ 带宽} & \textbf{环路主极点位置} & \textbf{稳定性评估} \\
\midrule
$0\,\mathrm{dB}$ ($1\times$) & $38\,\mathrm{MHz}$ & $10\,\mathrm{MHz}$ & 稳定 (裕度受限) \\
$20\,\mathrm{dB}$ ($10\times$) & $8\,\mathrm{MHz}$ & $2\,\mathrm{MHz}$ & 极佳 \\
$40\,\mathrm{dB}$ ($100\times$) & $700\,\mathrm{kHz}$ & $200\,\mathrm{kHz}$ & 优良 (契合系统带宽) \\
$60\,\mathrm{dB}$ ($1000\times$) & $87\,\mathrm{kHz}$ & $30\,\mathrm{kHz}$ & 稳定 (带宽严重收缩) \\
\bottomrule
\end{tabular}
\end{table}

\textbf{结果分析：} 根据“增益-带宽积守恒”定律，随着需求增益的提升，系统带宽不可避免地线性下降。如表\ref{tab:fda_stb}所示，当闭环增益设定为 $40\,\mathrm{dB}$ 时，系统带宽恰好在 $700\,\mathrm{kHz}$ 附近，完美契合了本项目的数据采集带宽需求。因此，在物理架构规划中，将该 FDA 的实际级联增益锁定在 $20\,\mathrm{dB} \sim 40\,\mathrm{dB}$ 之间，不仅能保证绝对的环路稳定性，还能确保高频信号不被过度衰减。

\subsection{共模抑制比（CMRR）与电源抑制比（PSRR）}

\textbf{仿真设置：}采用交流（AC）仿真机制。
1. **CMRR测试：** 在 FDA 的两个输入端同时施加极性相同的交流共模信号，测得共模增益（$A_{cm}$）；
2. **PSRR测试：** 在 VDD 供电轨叠加交流扰动信号，测得电源到输出的增益（$A_{vdd}$）。
将两者分别与预先测得的差模增益（$A_{dm}$）取对数比值，计算出 CMRR 和 PSRR。
仿真所用 testbench 如图\ref{fig:fda_cmrr_psrr_tb}所示。

\begin{figure}[H]
\centering
\subfloat[]{
  \label{fig:fda_cmrr_tb}
  \includegraphics[width=0.45\linewidth]{fda_cmrr_tb.png}
}
\subfloat[]{
  \label{fig:fda_psrr_tb}
  \includegraphics[width=0.45\linewidth]{fda_psrr_tb.png}
}
\caption{(a)共模抑制比（CMRR）仿真testbench；(b)电源抑制比（PSRR）仿真testbench（以正电源为例）}
\label{fig:fda_cmrr_psrr_tb}
\end{figure}

\textbf{仿真输出：} CMRR 与 PSRR 随频率变化的谱线。
四个不同工艺角下的 CMRR/PSRR 频率响应如图\ref{fig:fda_cmrr_psrr_curve}所示。

\begin{figure}[H]
    \centering
    \includegraphics[width=0.6\linewidth]{fda_cmrr_psrr_curve.png}
    \caption{FDA共模抑制比（CMRR）与电源抑制比（PSRR）仿真结果}
    \label{fig:fda_cmrr_psrr_curve}
\end{figure}

\textbf{仿真结果：} 在系统核心带宽内的抑制能力如表\ref{tab:fda_cmrr_psrr}所示。

\begin{table}[H]
\centering
\caption{FDA 的 CMRR 与 PSRR 仿真指标}
\label{tab:fda_cmrr_psrr}
\begin{tabular}{ccc}
\toprule
\textbf{抗扰指标} & \textbf{低频极限值} & \textbf{在 $1\,\mathrm{MHz}$ 有效带宽内的最差值} \\
\midrule
共模抑制比 (CMRR) & $> 280\,\mathrm{dB}$ & $\approx 280\,\mathrm{dB}$ \\
电源抑制比 (PSRR) & $> 250\,\mathrm{dB}$ & $\approx 250\,\mathrm{dB}$ \\
\bottomrule
\end{tabular}
\end{table}

\textbf{结果分析：} 在 14-bit 高精度混合信号系统中，由于后级 SAR ADC 存在频繁的开关动作，VDD 供电轨极易受到数字翻转噪声的污染；同时前端传感器环境易引入环境共模干扰。高达 $250\,\mathrm{dB}$ 的电源抑制比与 $280\,\mathrm{dB}$ 的共模抑制比（此为理想匹配状态下的极限值），充分证明了全差分对称架构的卓越抗干扰能力，能够将电源纹波与共模漂移彻底阻挡在信号链路之外。

\subsection{噪声分析}

\textbf{仿真设置：}将 FDA 配置为闭环工作模式（模拟其实际驱动后级的工作状态），输入共模偏置于 $1.65\,\mathrm{V}$。运行 Noise 仿真，噪声积分频率范围严谨设定为 $1\,\mathrm{Hz}$ 至 $10\,\mathrm{MHz}$，以覆盖所有的带内底噪及宽带热噪声混叠效应。
仿真所用 testbench 如图\ref{fig:fda_noise_tb}所示。

\begin{figure}[H]
    \centering
    \includegraphics[width=0.5\linewidth]{fda_noise_tb.png}
    \caption{FDA噪声分析仿真testbench}
    \label{fig:fda_noise_tb}
\end{figure}

\textbf{仿真输出：} FDA 折算到输入端的等效输入噪声电压谱密度（Input-Referred Noise Density）及其在频域上的累积积分曲线。
不同工艺角下噪声密度随频率变化的对数曲线如图\ref{fig:fda_noise_curve}所示。

\begin{figure}[H]
    \centering
  \includegraphics[width=0.6\linewidth]{fda_noise_curve.png}
  \caption{FDA等效输入噪声谱密度与积分噪声仿真结果}
  \label{fig:fda_noise_curve}
\end{figure}

\textbf{仿真结果：} 核心噪声指标如表\ref{tab:fda_noise}所示。

\begin{table}[H]
\centering
\caption{FDA 噪声特性仿真数据}
\label{tab:fda_noise}
\begin{tabular}{cc}
\toprule
\textbf{噪声评估项} & \textbf{仿真结果值} \\
\midrule
有效带内等效输入噪声密度 (@ $1\,\mathrm{MHz}$) & $\approx 2\,\mathrm{nV/\sqrt{Hz}}$ \\
折算至输入端的总积分噪声 ($0 \sim 1\,\mathrm{MHz}$) & $45\,\mathrm{\mu V_{rms}}$ \\
\bottomrule
\end{tabular}
\end{table}

\textbf{结果分析：} 在 14-bit（$\mathrm{LSB} \approx 183\,\mathrm{\mu V}$）的高分辨率系统中，FDA 的宽带热噪声和低频闪烁噪声（$1/f$ 噪声）是吞噬信噪比的元凶。从根本的物理设计角度，本设计通过为 FDA 第一级差分对管分配极大的尾电流（提升 $g_m$），并适度放大了输入管的栅面积（压低 $1/f$ 拐角），成功将带内噪声密度压制在极低的 $2\,\mathrm{nV/\sqrt{Hz}}$ 水平。最终 $45\,\mathrm{\mu V_{rms}}$ 的总积分噪声完美守住了系统的苛刻底噪红线。

\subsection{瞬态阶跃响应与压摆率}

\textbf{仿真设置：}将 FDA 配置为闭环模式，并在输出端挂载严苛的 $40\,\mathrm{pF}$ 大负载电容（等效模拟甚至超过了 SAR ADC 中 CDAC 阵列的极端采样电容）。在输入端施加幅度趋近满量程差分的陡峭阶跃信号。
仿真所用 testbench 如图\ref{fig:fda_tran_tb}所示。

\begin{figure}[H]
    \centering
    \includegraphics[width=0.5\linewidth]{fda_tran_tb.png}
    \caption{FDA瞬态阶跃响应与压摆率仿真testbench}
    \label{fig:fda_tran_tb}
\end{figure}

\textbf{仿真输出：} FDA 差分输出电压随时间建立的时域响应波形。重点量测大信号跳变斜率（Slew Rate, SR）以及小信号稳态建立至指定误差带内的所需时间（Settling Time）。
大信号阶跃响应及局部建立细节放大图如图\ref{fig:fda_tran_curve}所示。

\begin{figure}[H]
    \centering
  \includegraphics[width=0.8\linewidth]{fda_tran_curve.png}
  \caption{FDA瞬态阶跃响应及建立波形}
  \label{fig:fda_tran_curve}
\end{figure}

\textbf{仿真结果：} 动态大信号指标提取如表\ref{tab:fda_tran}所示。

\begin{table}[H]
\centering
\caption{FDA 大信号阶跃建立与压摆率参数}
\label{tab:fda_tran}
\begin{tabular}{cc}
\toprule
\textbf{动态响应指标} & \textbf{仿真结果值} \\
\midrule
差分压摆率 (Slew Rate) & $800\,\mathrm{V/\mu s}$ \\
建立时间 ($t_s$, 建立至 $500\,\mathrm{\mu V}$ 误差带) & $2.0\,\mathrm{\mu s}$ \\
\bottomrule
\end{tabular}
\end{table}

\textbf{结果分析：} 驱动沉重的容性负载并实现无振铃建立，是验证运放驱动能力的关键。得益于大偏置电流的设计，FDA 的内部充放电节点展现出了高达 $800\,\mathrm{V/\mu s}$ 的压摆率，这意味着其大信号跳变极为迅猛，非线性摆动时间被极度压缩。建立至 $500\,\mathrm{\mu V}$（近乎半个 LSB）仅耗时 $2\,\mathrm{\mu s}$，这为系统的采样时钟分配留足了时间窗口，完全规避了动态建立不足带来的谐波失真。

\subsection{总谐波失真（THD）}

\textbf{仿真设置：}在闭环模式下，向 FDA 注入频率为 $100\,\mathrm{kHz}$、幅度为满量程（$3\,\mathrm{V_{pp}}$ 差分）的纯净正弦波信号。运行高精度的瞬态仿真（Transient）以获取长时稳定波形，并利用 Virtuoso 内部的 Spectrum 工具对多个完整信号周期进行快速傅里叶变换（FFT），提取各次谐波能量。
仿真所用 testbench 如图\ref{fig:fda_thd_tb}所示。

\begin{figure}[H]
    \centering
    \includegraphics[width=0.5\linewidth]{fda_thd_tb.png}
    \caption{FDA总谐波失真（THD）仿真testbench}
    \label{fig:fda_thd_tb}
\end{figure}

\textbf{仿真输出：} FDA 满量程差分输出信号的 FFT 频谱图。图中清晰标定基波幅值与高次（主要为二次、三次）谐波分量的底噪水平。
FFT 频谱分析结果如图\ref{fig:fda_thd_curve}所示。

\begin{figure}[H]
  \centering
  \includegraphics[width=0.6\linewidth]{fda_thd_curve.png}
  \caption{FDA满量程输出FFT频谱与THD仿真结果}
  \label{fig:fda_thd_curve}
\end{figure}

\textbf{仿真结果：} 在不同的工艺边界条件下，提取出的满量程 THD 数据如表\ref{tab:fda_thd}所示。

\begin{table}[H]
\centering
\caption{FDA 在不同工艺角下的总谐波失真 (THD) 仿真结果}
\label{tab:fda_thd}
\begin{tabular}{ccc}
\toprule
\textbf{测试设定} & \textbf{工艺角 (Corners)} & \textbf{测得 THD (dB)} \\
\midrule
满量程, $100\,\mathrm{kHz}$ & Nominal (典型) & $-96.2\,\mathrm{dB}$ \\
满量程, $100\,\mathrm{kHz}$ & Asym\_FS & $-92.5\,\mathrm{dB}$ \\
满量程, $100\,\mathrm{kHz}$ & Asym\_SF & $-92.8\,\mathrm{dB}$ \\
满量程, $100\,\mathrm{kHz}$ & BC (最快) & $-98.1\,\mathrm{dB}$ \\
满量程, $100\,\mathrm{kHz}$ & WCLowtemp (最慢) & $-88.4\,\mathrm{dB}$ \\
\bottomrule
\end{tabular}
\end{table}
*(注：若增益过大导致信号超出供电轨满量程，管子进入线性区，THD 将严重恶化，不在可用线性范围内。)*

\textbf{结果分析：} 14-bit 高精度系统要求其模拟前端具有极其优秀的线性度（理论量化信噪比底线约 $-86\,\mathrm{dB}$）。全差分（FDA）架构在物理原理上自然抵消了绝大部分偶次谐波（如 HD2），使得系统的非线性主要受限于三次谐波（HD3）。仿真结果清晰展示，即便在最恶劣的 WCLowtemp 工艺角下，THD 依然优于 $-88\,\mathrm{dB}$，典型状态更是深达 $-96\,\mathrm{dB}$，这意味着 FDA 在进行满量程信号放大时，其非线性失真完全被掩埋在量化噪声之下。

\subsection{蒙特卡洛分析}

\textbf{仿真设置：}加载 TSMC $180\,\mathrm{nm}$ \texttt{mc} 统计工艺模型，设置仿真迭代次数 $N=100$ 次。综合考量工艺的随机偏差（Process Variation）与局部器件失配（Mismatch）。在闭环状态且差分输入端完全短接置零（$V_{id}=0\,\mathrm{V}$）的条件下，量测 FDA 输出端的直流失调电压（DC Offset Voltage）。
仿真所用 testbench 同前，如图\ref{fig:fda_thd_tb}所示。

\textbf{仿真输出：} 记录 100 次迭代的差分输出失调电压，并绘制统计散布直方图（Histogram），叠加高斯（Gaussian）拟合曲线。
FDA 输出失调电压的 Monte Carlo 统计分布如图\ref{fig:fda_mc_curve}所示。

\begin{figure}[H]
  \centering
  \includegraphics[width=0.6\linewidth]{fda_mc_curve}
  \caption{FDA差分输出失调电压蒙特卡洛仿真结果}
  \label{fig:fda_mc_curve}
\end{figure}

\textbf{仿真结果：} 基于数据提取，其统计分布特征参数如下表\ref{tab:fda_mc}所示：

\begin{table}[H]
\centering
\caption{FDA 输出直流失调电压蒙特卡洛统计参数}
\label{tab:fda_mc}
\begin{tabular}{ccc}
\toprule
\textbf{统计项目} & \textbf{平均值 ($\mu$)} & \textbf{标准差 ($1\sigma$)} \\
\midrule
输出直流失调电压 & $\approx 6.66\,\mathrm{aV}$ & $65.32\,\mathrm{aV}$ \\
\bottomrule
\end{tabular}
\end{table}

\textbf{结果分析：} 在基于理想原理图（Schematic-level）的蒙特卡洛仿真中，失调电压处于 $\mathrm{aV}$（阿伏特）级别，这一结果表明 FDA 的核心电路拓扑结构（包含电流镜与差分对）具有完美的本征对称性，不存在系统性失调漏洞。但在实际的物理流片中，由于版图布线的寄生不对称和掺杂浓度的随机波动，实际失调通常会落在 $\mathrm{mV}$ 量级。这深刻指导我们在后续版图（Layout）设计阶段，必须严格采用共源共栅交叉匹配（Common-Centroid）、添加外围伪管（Dummy Device）等高级版图技术，以最大程度捍卫这种原理图级别的对称性，避免失调电压挤占 ADC 的动态量程。

\subsection{实际指标总结}

基于上述严谨、全方位的交直流与瞬态仿真验证体系，本设计中全差分放大器（FDA）的核心物理性能指标最终收敛，并汇总对比于表\ref{tab:fda_final_specs}中：

\begin{table}[H]
\centering
\caption{FDA 模块实际核心仿真指标汇总}
\label{tab:fda_final_specs}
\begin{tabular}{lp{3.5cm}p{7cm}}
\toprule
\textbf{设计指标项} & \textbf{最终仿真达成值} & \textbf{工程价值与设计备注} \\
\midrule
开环直流增益 ($A_{v0}$) & $138\,\mathrm{dB}$ & 极大地压减了系统闭环稳态误差，确保高精度。 \\
闭环 $-3\,\mathrm{dB}$ 带宽 & $\approx 700\,\mathrm{kHz}$ (@ 40dB) & 完美契合前端数据采集系统的有效带宽需求。 \\
环路相位裕度 (PM) & $63^\circ$ & 巴特沃斯极点分布，保障阶跃响应无严重振铃。 \\
输入共模范围 (ICMR) & $0.4\,\mathrm{V} \sim 2.2\,\mathrm{V}$ & 单电源 $3.3\mathrm{V}$ 下提供宽阔的前级适应性。 \\
共模抑制比 (CMRR) & $> 280\,\mathrm{dB}$ & 有效过滤前端传感器的共模漂移干扰。 \\
电源抑制比 (PSRR) & $> 250\,\mathrm{dB}$ & 高度免疫 VDD 供电轨上的数字开关纹波。 \\
折算输入积分噪声 & $45\,\mathrm{\mu V_{rms}}$ & 绝对底噪低于 14-bit 系统的 LSB 误差红线。 \\
满量程总谐波失真 & $\le -88\,\mathrm{dB}$ (最恶劣) & 全差分架构滤除偶次谐波，保障系统极高线性度。 \\
大信号建立时间 & $2.0\,\mathrm{\mu s}$ (@ $40\mathrm{pF}$ 负载) & 赋予后级 SAR ADC 充足的采样建立追踪时间。 \\
\bottomrule
\end{tabular}
\end{table}


% -------------------------------
% 可编程增益放大器（PGA）单模块仿真
% @ 模块基本测试条件：
%   1. 电源3.3V
%   2. 标准共模输出1.65V
% @ 模块技术细节：
%   1. 采用全差分架构，利用R-2r网络实现可编程增益
%   2. Rg与Rf均采用8bit的r-2r网络，实现1/255-255倍间任意分数增益
% -------------------------------

\section{可编程增益放大器（PGA）单模块仿真}

\subsection{输出共模范围稳定性}

\textbf{仿真设置：}进行直流（DC）参数扫描仿真。
配置 PGA 供电电压为 $3.3\,\mathrm{V}$，CMFB 参考电压（即输出共模目标）精准设定为 $1.65\,\mathrm{V}$。在输入端同步扫描共模电平，扫描范围覆盖 $0.5\,\mathrm{V}$ 至 $2.5\,\mathrm{V}$，通过输出端内置的共模检测电阻网络提取实际输出共模电压。
仿真所用 testbench 如图 \ref{fig:pga_ocmr_tb} 所示。

\begin{figure}[H]
    \centering
    \includegraphics[width=0.5\linewidth]{pga_ocmr_tb.png}
    \caption{PGA 输出共模范围稳定性仿真 testbench}
    \label{fig:pga_ocmr_tb}
\end{figure}

\textbf{仿真输出：} PGA 实际输出共模电压随输入共模电压扫描变化的直流响应曲线。
PGA 输出共模电平随输入共模变化的特性曲线如图 \ref{fig:pga_ocmr_curve} 所示。

\begin{figure}[H]
  \centering
  \includegraphics[width=0.7\textwidth]{pga_ocmr_curve.png}
  \caption{PGA 输出共模电平稳定性仿真结果}
  \label{fig:pga_ocmr_curve}
\end{figure}

\textbf{仿真结果：} 当输入共模在 $0.5\,\mathrm{V} \sim 2.5\,\mathrm{V}$ 的宽泛区间内剧烈波动时，PGA 的输出共模电压被严格钳位在 $1.65\,\mathrm{V}$ 中心点附近，最大波动幅度 $< 0.01\,\mathrm{mV}$。

\textbf{结果分析：} 在多级级联的模拟前端（AFE）链路中，前级微弱信号在经过高增益放大后，其直流基线极易发生不规则的共模漂移（Baseline Wander）。仿真表明，得益于 PGA 内部具备极高低频环路增益的共模反馈（CMFB）网络，PGA 能够完美隔离并抑制前级的共模漂移，为后级 SAR ADC 的全差分采样阵列提供极其稳定、精准且抗干扰的中心共模电平保障。

\subsection{增益步进准确度}

\textbf{仿真设置：}进行交流（AC）频响仿真。在 PGA 差分输入端施加幅度为 $1\,\mathrm{mV}$ 的交流小信号。通过向内部寄存器写入数字逻辑控制字，改变 R-2R 反馈电阻网络的继电器/开关阵列接入状态。采用参数化扫描（Parametric Analysis）遍历从最小增益（控制字 \texttt{0x00}）到最大增益（控制字 \texttt{0xFF}）的各个核心增益档位。
仿真所用 testbench 如图 \ref{fig:pga_gain_step_tb} 所示。

\begin{figure}[H]
    \centering
    \includegraphics[width=0.5\linewidth]{pga_gain_step_tb.png}
    \caption{PGA 增益步进与控制逻辑仿真 testbench}
    \label{fig:pga_gain_step_tb}
\end{figure}

\textbf{仿真输出：} 提取多组不同控制字下，PGA 闭环增益低频平坦区的绝对增益数值，并计算其与理论分数阶增益的相对误差百分比。
PGA 在典型增益配置下的幅频响应簇如图\ref{fig:pga_gain_step_curve}所示。

\begin{figure}[H]
  \centering
  \includegraphics[width=0.5\linewidth]{pga_gain_step_curve.png}
  \caption{PGA 闭环增益步进扫描仿真结果}
  \label{fig:pga_gain_step_curve}
\end{figure}

\textbf{仿真结果：} 增益调节范围成功覆盖了设计的 $1/255 \sim 255\,\mathrm{V/V}$ 动态区间，步进增量体现出完美的单调性，且各档位实际增益与 R-2R 理论计算值高度吻合，全区间最大增益相对误差被严格控制在 $\le 0.01\%$ 以内。

\begin{table}[H]
\centering
\caption{PGA 增益步进准确度与单调性测试结果}
\label{tab:pga_gain_error}
\begin{tabular}{cccc}
\toprule
\textbf{目标增益 (V/V)} & \textbf{实际仿真增益 (V/V)} & \textbf{增益误差 (\%)} & \textbf{是否单调} \\
\midrule
$1/16$ & $0.062503$ & $+0.0048\%$ & 是 \\
$1/8$  & $0.125006$ & $+0.0048\%$ & 是 \\
$1/4$  & $0.249980$ & $-0.0080\%$ & 是 \\
$1/2$  & $0.499970$ & $-0.0060\%$ & 是 \\
$1$    & $1.000020$ & $+0.0020\%$ & 是 \\
$2$    & $1.999950$ & $-0.0025\%$ & 是 \\
$4$    & $3.999800$ & $-0.0050\%$ & 是 \\
$8$    & $8.000400$ & $+0.0050\%$ & 是 \\
$16$   & $15.99900$ & $-0.0062\%$ & 是 \\
\bottomrule
\end{tabular}
\end{table}

\textbf{结果分析：} PGA 的增益步进准确度直接决定了工业宽动态范围信号在被归一化处理后的系统整体线性度。采用高精度的 R-2R 网络实现了精细的分数阶增益调节，仿真证实了在充分考虑开关导通电阻（$R_{\mathrm{on}}$）影响的条件下，优化的闭环电阻网络并未引起严重的直流增益衰减。全档位误差低于 $0.01\%$ 表明其静态增益匹配度完全满足了 14 位系统对增益误差（Gain Error）的苛刻容忍极限。

\subsection{压摆率与建立时间}

\textbf{仿真设置：}进行瞬态大信号（Transient）仿真。在 PGA 的全差分输出端挂载高达 $40\,\mathrm{pF}$ 的差分等效电容负载（模拟 SAR ADC 内部 CDAC 阵列在采样相位的极端等效电容）。将 PGA 锁定于最大增益配置下，在输入端注入一阶跃激励，促使输出产生接近满量程的急剧电压跳变。
仿真所用 testbench 如图 \ref{fig:pga_tran_tb} 所示。

\begin{figure}[H]
    \centering
    \includegraphics[width=0.5\linewidth]{pga_tran_tb.png}
    \caption{PGA 大信号压摆率与建立时间仿真 testbench}
    \label{fig:pga_tran_tb}
\end{figure}

\textbf{仿真输出：} 捕捉 PGA 差分输出电压随时间建立的精确时域波形，提取其上升/下降沿的最大斜率（压摆率 Slew Rate），并局部极度放大以观测信号平稳切入目标精度误差带的时序细节。
PGA 满量程阶跃响应及 14-bit 精度建立细节如图\ref{fig:pga_tran_curve}所示。

\begin{figure}[H]
  \centering
  \includegraphics[width=0.7\textwidth]{pga_tran_curve.png}
  \caption{PGA 大信号瞬态响应与建立时间仿真结果}
  \label{fig:pga_tran_curve}
\end{figure}

\textbf{仿真结果：} 差分压摆率峰值达到 $200\,\mathrm{V/\mu s}$，且输出信号平滑无显著欠阻尼过冲。最终信号稳定建立至 $\pm 0.5\,\mathrm{LSB}$（针对 14 位系统，即约 $0.003\%$ 误差带）所需的总建立时间为 $t_s = 800\,\mathrm{ns}$。

\textbf{结果分析：} 稳定驱动沉重的 CDAC 容性负载并实现极速、高精度的无损建立，是 PGA 模块设计中最具挑战性的任务。高达 $200\,\mathrm{V/\mu s}$ 的压摆率证明了 PGA 的输出级偏置架构具备强大的瞬态充放电能力。对于 $1\,\mathrm{MS/s}$ 采样率的系统（采样周期 $1\,\mathrm{\mu s}$），$800\,\mathrm{ns}$ 的仿真建立时间虽然紧凑，但在分配了非对称时钟相位的跟踪期（Tracking Phase）内，依然能够完美保障后级 ADC 对满量程信号的绝对高精度截取。

\subsection{噪声分析}

\textbf{仿真设置：}运行高精度 Noise 仿真。为剥离最恶劣的噪声贡献基线，将 PGA 配置于最大增益状态（此状态下运放本体噪声折算到输入端最显著）。积分频带严格设定为 $1\,\mathrm{Hz}$ 至 $10\,\mathrm{MHz}$，提取 PGA 的等效输入噪声（Input-Referred Noise, IRN）。
仿真所用 testbench 与增益步进测试相同，如图 \ref{fig:pga_gain_step_tb} 所示。

\textbf{仿真输出：} PGA 等效输入噪声电压谱密度（$\mathrm{nV/\sqrt{Hz}}$）及其在全有效频段内的累积积分噪声电压轨迹。
PGA 等效输入噪声特性曲线如图\ref{fig:pga_noise_curve}所示。

\begin{figure}[H]
  \centering
  \includegraphics[width=0.5\linewidth]{pga_noise_curve.png}
  \caption{PGA 等效输入噪声谱密度仿真结果}
  \label{fig:pga_noise_curve}
\end{figure}

\textbf{仿真结果：} 在输入频率为 $100\,\mathrm{kHz}$ 时，测得 PGA 的等效输入噪声密度为 $5\,\mathrm{nV/\sqrt{Hz}}$，在整个系统带宽内的总积分噪声电压为 $67\,\mathrm{\mu V_{rms}}$。

\textbf{结果分析：} 根据 Friis 噪声级联公式（$F_{\mathrm{total}} = F_1 + \frac{F_2 - 1}{G_1} + \dots$），在多级级联架构中，后级电路的等效输入噪声折算到系统最前端时，会被前向增益强烈衰减。由于本系统前级 TIA 提供了高达 $100\,\mathrm{dB\Omega}$（$10^5\,\mathrm{V/A}$）的巨大跨阻增益，PGA 自身贡献的 $67\,\mathrm{\mu V_{rms}}$ 积分噪声折算至传感器输入端等效仅为皮安（$\mathrm{pA}$）级别，完全被淹没在第一级的本底噪声之下。仿真确认 PGA 的内部电路设计足够纯净，未引入异常的高频噪声峰化（Noise Peaking），不会对整个 AFE 链路的核心噪声预算构成威胁。

\subsection{实际指标总结}

基于对可编程增益放大器（PGA）全维度的参数扫描与时域/频域协同仿真验证，本模块的各项核心工程物理指标已被精准提取，并汇总定型于表\ref{tab:pga_final_specs}中：

\begin{table}[H]
\centering
\caption{PGA 模块实际核心仿真指标汇总}
\label{tab:pga_final_specs}
\begin{tabular}{lp{3.5cm}p{7cm}}
\toprule
\textbf{设计关键指标项} & \textbf{最终仿真达成量测值} & \textbf{工程应用价值与设计备注} \\
\midrule
系统推荐增益范围 & $1/255 \sim 255\,\mathrm{V/V}$ & 支持极宽范围的分数阶动态调节，满足复杂场景信号归一化。 \\
静态中心输出共模 & $1.65\,\mathrm{V} \pm 0.01\,\mathrm{mV}$ & 依托高增益 CMFB 环路，提供近乎绝对死锁的共模基线。 \\
全局最大增益误差 & $\le 0.01\%$ & 杜绝了大范围增益切换时的非线性跳变与失真。 \\
差分信号压摆率 & $200\,\mathrm{V/\mu s}$ (@ $40\,\mathrm{pF}$) & 输出级具备超强充放电能力，无惧沉重开关电容冲击。 \\
大信号阶跃建立时间 & $800\,\mathrm{ns}$ ($\pm 0.5\,\mathrm{LSB}$) & 建立曲线平滑，完美匹配 $1\,\mathrm{MS/s}$ 高速采样系统的追踪相。 \\
模块等效积分噪声 & $67\,\mathrm{\mu V_{rms}}$ (@ 最大增益) & 极低的热噪声底色，折算至系统前端后几乎完全可忽略。 \\
\bottomrule
\end{tabular}
\end{table}

% -------------------------------
% 比较器（Comp）单模块仿真
% -------------------------------

\section{比较器（Comp）单模块仿真}


% -------------------------------
% 电容数字模拟转换器（CDAC）单模块仿真
% @ 模块基本测试条件：
%   1. 电源3.3V
%   2. 输入差分信号共模范围为1.65V，差分电压信号输入为1-2V
% @ 模块技术细节：
%   1. 采用14位电容数字模拟转换器，由两组7位电容数字模拟转换器组成，中间采用桥接电容Cc连接
%   2. 单位电容值为50fF
% -------------------------------



